% The user playbook, in English. The look is in sheriff-playbook.sty, shared
% with the Spanish version (sheriff-playbook-es.tex): a change to what the user
% types or reads goes into both in the same change. scripts/playbook.sh builds
% both.
\documentclass[11pt,a4paper]{article}
\usepackage[english]{sheriff-playbook}

\begin{document}

\makecover{Sheriff}{Kaizten's code standards\\in Claude Code, Codex and Antigravity}{Installation, daily use, continuous integration and troubleshooting.}

\section*{How to use this playbook}

This playbook is for anyone who writes code with Claude Code, Codex or Antigravity and wants
that code to meet Kaizten's standards from the first line. You do not need to
know how the tools are built: this is what to install, what to ask for and what
to expect at each point.

Read sections 1 to 3 once, when you start. From section 4 on, the playbook is
organised in \emph{plays}: each one is a concrete situation, with when to use
it, what to ask, what happens and how to tell it is done. The last page is a
reference card to keep at hand.

\begin{tip}
\textbf{If you only have five minutes.}
\begin{enumerate}[topsep=2pt,leftmargin=1.3em]
  \item Run the install line; no GitHub account is needed. If something is
        missing, such as Java 17 or Docker, the installer tells you
        (section~\ref{sec:install}). On Linux and macOS:
\end{enumerate}
\begin{command}[basicstyle=\ttfamily\scriptsize]{bash}
curl -fsSL https://github.com/kaizten/sheriff-mcp-server/releases/latest/download/install.sh \
  | bash
\end{command}
\noindent\hspace*{1.3em}On Windows, in PowerShell:
\begin{command}[basicstyle=\ttfamily\footnotesize]{PowerShell}
$u = 'https://github.com/kaizten/sheriff-mcp-server/releases/latest/download';
& ([scriptblock]::Create((irm "$u/install.ps1")))
\end{command}
\begin{enumerate}[resume,leftmargin=1.3em]
  \item Open Claude Code, Codex or Antigravity in your project and ask
        \ask{check the code standards}.
\end{enumerate}
\end{tip}

\vspace{0.2\baselineskip}
{\sffamily\bfseries\Large\color{pbbrand}Contents}\par
\vspace{-0.7\baselineskip}
{\color{pbrule}\rule{\linewidth}{0.8pt}}
\makeatletter\@starttoc{toc}\makeatother

\clearpage
% ---------------------------------------------------------------------------
\section{What Sheriff is and what it does for you}
\label{sec:what}

\textbf{Sheriff} is Kaizten's code standards analyzer. It checks, for example,
that every method has its JavaDoc, that imports are sorted, that every
\code{if} has braces, or that the code respects the project's architecture
(hexagonal, DDD). Its rules are compiled into a Docker image,
\code{kaizten/sheriff}: they are not a style guide to interpret, but checks that
pass or fail.

These tools bring Sheriff to where you write code. With them, your assistant:

\begin{itemize}
  \item checks the component \textbf{before} touching it, and clears the errors
        that were already there;
  \item reads the rules before writing, so as not to introduce new errors;
  \item applies Sheriff's own \textbf{automatic repairs} first, which spend no
        tokens, and fixes by hand only what is left;
  \item does not call the work done while errors remain or the project's tests
        fail.
\end{itemize}

\begin{center}
\begin{tikzpicture}[node distance=4mm and 7mm]
  \node[pbstrong,minimum width=17mm,minimum height=13mm] (you) {You};
  \node[pbbox,right=of you,minimum width=29mm,minimum height=13mm] (assistant) {\textbf{Your assistant}\\[1pt]\scriptsize Claude Code, Codex\\\scriptsize or Antigravity};
  \node[pbbox,right=15mm of assistant,minimum width=37mm] (server) {\textbf{MCP server}\\[1pt]\scriptsize five tools};
  \node[pbbox,below=of server,minimum width=37mm] (hooks) {\textbf{Hooks}\\[1pt]\scriptsize before an edit, at the end};
  \node[pbbox,below=of hooks,minimum width=37mm] (agent) {\textbf{Repair agent}\\[1pt]\scriptsize \code{sheriff_autofix}, \code{--agent}};
  \node[pbbox,below=of agent,minimum width=37mm] (check) {\textbf{CI check}\\[1pt]\scriptsize \code{--check}};
  \begin{scope}[on background layer]
    \node[draw=pbaccent,dashed,line width=0.9pt,fill=pbaccentsoft,rounded corners=5pt,inner sep=3.5mm,fit=(server)(hooks)(agent)(check)] (jar) {};
  \end{scope}
  \node[anchor=south,font=\sffamily\scriptsize\bfseries,text=pbaccent] at (jar.north) {sheriff-mcp.jar};
  \node[pbstrong,right=11mm of jar,minimum width=23mm,minimum height=15mm] (sheriff) {Sheriff\\[1pt]\scriptsize\color{pbaccentsoft}Docker image};
  \node[pbsoft,minimum width=29mm] (ci) at (assistant |- check) {CI, terminal\\Maven};
  \draw[pbarrow] (you) -- (assistant);
  \draw[pbarrow] (assistant) -- node[pblabel,above]{MCP} (server);
  \draw[pbarrow,draw=pbmuted] (hooks.west) -| node[pblabel,pos=0.3,below,align=center]{stop it while\\errors remain} (assistant.south);
  \draw[pbarrow,draw=pbmuted] (ci) -- (check);
  \draw[pbarrow] (jar.east) -- (sheriff.west);
\end{tikzpicture}
\end{center}

Everything lives in one file, \code{sheriff-mcp.jar}, installed once per
machine. Your project copies nothing and needs no configuration: the tools work
out what to analyze, and with which rules, from the folder you work in.

\subsection{What changes in your daily work}

\begin{itemize}
  \item \textbf{You ask for changes as always.} The assistant checks the
        standards on its own; you do not have to ask.
  \item \textbf{If the component already had errors, it fixes them first},
        even if you did not introduce them. That is on purpose: the code
        delivered is clean as a whole. When you do not want it, say so in the
        request: \ask{leave the errors that were already there}.
  \item \textbf{Before finishing, it checks again and runs the tests} of the
        project, and tells you which errors were there before it started.
  \item \textbf{If it tries to finish with errors} in what it changed, a hook
        sends the turn back. You do not have to watch for it.
\end{itemize}

\subsection{Six words worth knowing}

\begin{center}
\rowcolors{2}{pbcode}{white}
\begin{tabularx}{\linewidth}{@{\hspace{6pt}}K{33mm}L}
\headrow \thead{Word} & \thead{What it means here} \\
Component & The folder Sheriff analyzes: one module of the project, the one
            with \code{src/}. In a single-module project, the project
            itself. \\
Profile & The set of rules applied: \code{JAVA}, \code{JAVA_HEXAGONAL},
          \code{TYPESCRIPT}\ldots{} Worked out by itself. \\
Error & A breach of one rule, with its file, its line and its code (for
        example, \code{SortedImport}). \\
Automatic repair & One of Sheriff's \emph{fixers}: it repairs a rule with no AI
        model and no tokens spent. \\
Hook & A command Claude Code or Codex run by themselves at a set moment:
       before editing a file, or when a turn ends. \\
Task & Every analysis is recorded with an id (\code{Task 4047f176.}), useful
       to look something up or report it. \\
\end{tabularx}
\end{center}

% ---------------------------------------------------------------------------
\section{Installation}
\label{sec:install}

One line installs the latest release, on Linux, macOS and Windows. \textbf{The
same line updates} when a new one comes out. The installer checks what it
needs and, when something is missing, tells you what to do.

\subsection{The only thing you need first}
\label{sec:requirements}

Nothing: no GitHub account, and none of its tools. The install line downloads
the script from the repository's public releases page, and the script
downloads the rest from the same place. Java and Docker need not be there
first either: the installer looks for them and tells you what to do.

\subsection{The install line}

\needspace{4\baselineskip}
On Linux and macOS:

\begin{command}[basicstyle=\ttfamily\scriptsize]{bash}
curl -fsSL https://github.com/kaizten/sheriff-mcp-server/releases/latest/download/install.sh | bash
\end{command}

Without \code{curl}, \code{wget -qO-} with the same address does the same. On
Windows, in PowerShell (Windows PowerShell or PowerShell 7):

\begin{command}[basicstyle=\ttfamily\footnotesize]{PowerShell}
$u = 'https://github.com/kaizten/sheriff-mcp-server/releases/latest/download';
& ([scriptblock]::Create((irm "$u/install.ps1")))
\end{command}

\subsection{If you would rather read the script first}

The install line hands the script straight to \code{bash}. If you would rather
download it, read it and run it afterwards, do it this way; \code{--release}
tells it to install the release even when you run it inside a repository:

\begin{command}[basicstyle=\ttfamily\scriptsize]{bash}
curl -fsSLO https://github.com/kaizten/sheriff-mcp-server/releases/latest/download/install.sh
bash install.sh --release
\end{command}

\begin{command}[basicstyle=\ttfamily\footnotesize]{PowerShell}
$u = 'https://github.com/kaizten/sheriff-mcp-server/releases/latest/download';
irm "$u/install.ps1" -OutFile install.ps1
.\install.ps1 -Release
\end{command}

If PowerShell refuses to run the file, allow scripts in that window only with
\path{Set-ExecutionPolicy -Scope Process -ExecutionPolicy Bypass}. Both scripts
can also be downloaded from the repository's releases page on GitHub.

\subsection{What the installer checks}

\begin{center}
\rowcolors{2}{pbcode}{white}
\begin{tabularx}{\linewidth}{@{\hspace{6pt}}K{33mm}L}
\headrow \thead{Checks} & \thead{When it is missing} \\
Java 17 or newer & It looks in \code{JAVA_HOME}, on the \code{PATH} and in
  the folders JDKs are installed in, even when it is not on the \code{PATH},
  and sets the assistants up with its full path. With none, on Windows it
  installs Temurin 21 with \code{winget}; on Linux and macOS it stops and tells
  you where to get it (\path{https://adoptium.net}). \\
Docker & It goes on without it, warning that nothing can be analyzed until you
  have it. With \code{--install-docker} it installs it (on Linux with Docker's
  script or, on Arch, \code{pacman}, and adds you to the \code{docker} group;
  on macOS and Windows, Docker Desktop) and starts it when it is stopped. \\
Claude Code & If it is there, it registers the MCP server. If not, it prints the
  command to register it once you install it. \\
Codex & If it is there, it sets it up (section~\ref{sec:codex}). If not, it
  prints the command for when you install it. \\
Antigravity & If \code{agy} is there, it sets it up (section~\ref{sec:antigravity}).
  If not, it prints the command for when you install it. \\
The MCP server & At the end it starts it once, as your assistant will, and
  checks that it answers: \code{OK}, or \code{FAIL} with the cause. \\
\end{tabularx}
\end{center}

It writes Claude Code's hooks even when \code{claude} is not on the
\code{PATH}, as long as the folder \path{~/.claude} exists: they are only a
settings file, which the desktop app and the editor extensions read too. It
touches no project. Sheriff's rules (the \emph{catalog}) are not in the
download: the tools extract them from the image the first time they need them,
in about 6 seconds.

\begin{warning}
Docker Desktop asks you to accept its terms the first time it starts, and the
installer does not accept them for you. In a company of more than 250 people or
more than \$10 million in annual revenue they require a paid subscription.
\end{warning}

\subsection{What you will see}

This is the output of a real install on a machine with Claude Code, Codex and
Antigravity.
Between the third line and the fourth, the download of Sheriff's image shows
too: about 4\,GB the first time and, after that, only what changed.

\begin{answer}[Installer output]
Using Java 17 at /usr/bin/java
Downloading sheriff-mcp.jar from the latest release of kaizten/sheriff-mcp-server...
Installed sheriff-mcp 1.0.2 into /home/ana/.local/share/sheriff-agent
Sheriff's hooks are in /home/ana/.claude/settings.json, for every project Claude Code opens:
  before an edit     refuse new code on a component that already fails the standards
  when a turn ends   refuse to finish while what the turn changed has errors
  when a turn starts note what was already changed, so that it is not the turn's
sheriff_test, sheriff_fix, sheriff_guidelines and sheriff_task run without asking.
Where there is nothing Sheriff analyzes, they let everything through.
Undo with: java -jar sheriff-mcp.jar --uninstall-hooks --user
Added stdio MCP server sheriff with command: /usr/bin/java -jar /home/ana/.local/share/sheriff-agent/sheriff-mcp.jar to user config
File modified: /home/ana/.claude.json
Registered the MCP server as 'sheriff'. Restart open Claude Code sessions to pick it up.
Codex is set up for every project it opens:
  MCP server     'sheriff' in /home/ana/.codex/config.toml, with 600 s per call and sheriff_fix approved
  order of work  in /home/ana/.codex/AGENTS.md
Undo with: java -jar sheriff-mcp.jar --uninstall-codex
Sheriff's hooks are in /home/ana/.codex/hooks.json, for every project Codex opens:
  when a turn ends   refuse to finish while what the turn changed has errors
  when a turn starts note what was already changed, so that it is not the turn's
Codex runs them once you approve them: it asks the first time a session starts, or use /hooks.
Until then a turn can end with errors left: approve them in your first session.
Restart open Codex sessions to pick it up; the first one asks you to review the new hooks.
Antigravity is set up for every project agy opens:
  MCP server     'sheriff' in /home/ana/.gemini/config/mcp_config.json
  permissions    sheriff_test, sheriff_fix, sheriff_guidelines and sheriff_task run without asking,
                 in /home/ana/.gemini/antigravity-cli/settings.json
It has no Sheriff hooks: a turn there can end with errors left.
Undo with: java -jar sheriff-mcp.jar --uninstall-antigravity
Restart open agy sessions to pick it up.
The MCP server starts: OK, it answers as sheriff.

Sheriff is set up for Claude Code, Codex and Antigravity.
\end{answer}

It says which Java it will use, which version it installed and where and, for each assistant, which
hooks it turned on, which tools may run without asking, and how to undo each
one; then, that the server starts, and in its last line, which assistants it
set up. It sets up only the ones the
machine already has, and never installs one: an assistant that is not there
gets a line saying how to set it up once it is. As it asks,
\textbf{restart the Claude Code and Codex sessions you had open}. To check that
your assistants see the server: in Claude Code, \code{claude mcp list} must
show \code{sheriff:}~\ldots~\textcolor{pbtip}{\faCheck}~\code{Connected} (inside
a session, \code{/mcp}); in Codex, \code{codex mcp list} must show
\code{sheriff} as \code{enabled}, and in Antigravity, \code{agy mcp list}, the
same.

\subsection{Options}

Options go at the end of the install line: in bash, after \code{bash -s --};
in PowerShell, after the script block. For example:

\begin{command}{bash}
curl -fsSL .../install.sh | bash -s -- --pull-always
\end{command}
\begin{command}{PowerShell}
& ([scriptblock]::Create((irm "$u/install.ps1"))) -PullAlways
\end{command}

\begin{center}
\rowcolors{2}{pbcode}{white}
\begin{tabularx}{\linewidth}{@{\hspace{6pt}}>{\ttfamily\small}p{31mm}>{\ttfamily\small}p{26mm}L}
\headrow \thead{bash} & \thead{PowerShell} & \thead{What it does} \\
-{}-install-docker & -InstallDocker & Installs Docker when it is missing, and
  starts it when it is stopped. \\
 & -DockerWsl & Windows: uses the Docker installed inside WSL, with no Docker
  Desktop. Restart Claude Code afterwards. \\
-{}-pull-always & -PullAlways & Pulls a newer Sheriff image each time the server
  starts. Without it a newer image is only reported, once per session, so that
  rules do not change in the middle of someone's work. Give it again on every
  update. \\
-{}-no-pull & -NoPull & Does not pull the image now; the server pulls it when
  it needs it. \\
-{}-no-hooks & -NoHooks & Installs the MCP server without the hooks. \\
-{}-no-codex & -NoCodex & Leaves Codex's configuration alone. \\
-{}-no-antigravity & -NoAntigravity & Leaves Antigravity's configuration alone. \\
-{}-fail-fast & -FailFast & The hooks stop Sheriff at the first error: quicker,
  but they report only one. \\
\end{tabularx}
\end{center}

% ---------------------------------------------------------------------------
\section{Your first session}
\label{sec:first}

This is a real example, from start to finish, on a very small Maven project
called \code{basket}. Sheriff's answers are copied as they came, cut only where
\code{[...]} shows.

\subsubsection{1. Open the assistant in the project}
Open a terminal at the project's root (the folder with the \code{pom.xml}, the
\code{package.json} or the \code{src/}) and start \code{claude}, \code{codex} or
\code{agy}: Sheriff's answers and the steps are the same. The editor or
the desktop app work the same, as long as the folder open is the project's.

\subsubsection{2. Ask for the check}
Type \ask{check the code standards}. The assistant calls
\tool{sheriff_test}, which answers with a count:

\begin{answer}
basket has 4 error(s) under JAVA.
Paths below are under basket/src/main/java/demo/
  Basket.java: 4

By rule: JavaJavaDocCommentInMethodMessage6 2, JavaBracesIfStatementsMessage2 1, SortedImport 1

Task 4047f176.
Next step: call the sheriff_fix tool with component 'basket'. It applies every repair Sheriff has for these errors and lists what it cannot repair.
\end{answer}

\begin{note}
\textbf{How to read an answer.} The first line says which component was
analyzed, how many errors it has and under which profile. Then come the errors
per file and per rule, the task's id and, always on the last line,
\textbf{\code{Next step:}}, what the assistant does next. That step is worked
out from the state of the code, so it is not yours to decide. The assistant
sums the answer up for you in your own words.
\end{note}

\subsubsection{3. The automatic repairs}
The assistant follows the step and calls \tool{sheriff_fix}. Sheriff sorts the
import by itself and lists what is left, with the code it accepts for each
rule:

\begin{answer}
Ran every fixer Sheriff has for these errors on basket: 4 error(s) before, 3 after (1 repaired).
Files the fixers changed (1): Basket.java, under basket/src/main/java/demo/
Still outstanding: 3 error(s) in 1 file(s). This answer lists the 3 error(s) of the first 1 file(s), and the work is those files now:
Paths below are under basket/src/main/java/demo/

Basket.java
  - Braces are not used in if statement 'if (item == null)
    return;' (lines 10-10).  [JavaBracesIfStatementsMessage2]
  - Method 'add' in file 'Basket.java' does not have JavaDoc comment  [JavaJavaDocCommentInMethodMessage6]
  - Method 'size' in file 'Basket.java' does not have JavaDoc comment  [JavaJavaDocCommentInMethodMessage6]

What Sheriff accepts, for the rules above:
Every branch and loop body has braces, and its content starts on a line of its own:
    if (page < FIRST_PAGE) {
        throw new IllegalArgumentException(UNKNOWN);
    }
[...]
Task a3f3e1be.
Next step: fix by hand every error listed above, all of them, in basket/src/main/java/demo/Basket.java, each file from the last line up [...] Then call the sheriff_fix tool with component 'basket' again.
\end{answer}

\subsubsection{4. What is left, by hand}
The assistant adds the braces and the two JavaDocs, from the last line up (so
the line numbers do not move under it), and calls \tool{sheriff_fix} again,
which checks and repairs whatever those edits may have broken:

\begin{answer}
Ran every fixer Sheriff has for these errors on basket: 0 error(s) before, 0 after (0 repaired).

Task 50828b9f.
Next step: run the project's tests (cd '/home/ana/projects/basket' && mvn test). If they pass, the component is done. If you change any file to make them pass, call the sheriff_test tool with component 'basket' again.
\end{answer}

\subsubsection{5. The tests, and done}
The assistant runs the project's tests. If they pass, the component is done:
\textbf{0 errors and green tests}. It tells you so in those terms, and which
errors were there before it started.

\subsection{The work cycle}

Every request that changes code follows the same cycle. You do not have to ask
for it: the server tells the assistant when it connects, and every answer gives
it the next step.

\begin{center}
\begin{tikzpicture}[node distance=14mm and 6mm]
  \node[pbstrong,minimum width=25mm,minimum height=14mm] (ask) {Your request};
  \node[pbbox,right=of ask,minimum width=31mm,minimum height=14mm] (test) {\textbf{1 · Check}\\[1pt]\scriptsize\code{sheriff_test}};
  \node[pbbox,right=of test,minimum width=36mm,minimum height=14mm] (clean) {\textbf{2 · Clear what was there}\\[1pt]\scriptsize\code{sheriff_fix}, then by hand};
  \node[pbbox,below=of clean,minimum width=36mm,minimum height=14mm] (write) {\textbf{3 · Write the change}\\[1pt]\scriptsize with \code{sheriff_guidelines}};
  \node[pbbox,left=of write,minimum width=31mm,minimum height=14mm] (again) {\textbf{4 · Check again}\\[1pt]\scriptsize down to 0 errors};
  \node[pbbox,left=of again,minimum width=25mm,minimum height=14mm] (tests) {\textbf{5 · Tests}\\[1pt]\scriptsize the project's};
  \node[pbgold,below=9mm of tests,minimum width=25mm] (done) {\textbf{Done}\\[1pt]\scriptsize 0 errors, green tests};
  \draw[pbarrow] (ask) -- (test);
  \draw[pbarrow] (test) -- (clean);
  \draw[pbarrow] (clean) -- node[pblabel,right]{0 errors} (write);
  \draw[pbarrow] (write) -- (again);
  \draw[pbarrow] (again) -- (tests);
  \draw[pbarrow] (tests) -- (done);
  \path (clean) edge[pbarrow,draw=pbaccent,loop above,min distance=9mm,in=60,out=120] node[pblabel,above]{down to 0 errors} (clean);
  \draw[pbarrow,draw=pbaccent] ([xshift=-4mm]tests.south east) to[out=-60,in=-120,looseness=1.1] node[pblabel,below]{changed something? again} (again.south);
\end{tikzpicture}
\end{center}

% ---------------------------------------------------------------------------
\clearpage
\section{Everyday plays}
\label{sec:plays}

\begin{play}{Add a feature or fix a bug}
\when Any code change in a component Sheriff analyzes.
\asks Whatever you need, as always: \ask{add a method to Basket that returns
how many different items it holds}.
\happens
\begin{enumerate}
  \item The assistant checks the component and, if it has errors, clears them
        first (section~\ref{sec:first}).
  \item It reads the rules that bear on what it is about to write, and writes
        the change.
  \item It checks again and repairs whatever the change broke.
  \item It runs the project's tests.
  \item In Java, if it added a non-private method that no test calls, it writes
        that test, edge cases included (an empty list, a \code{null}).
\end{enumerate}
\done It reports 0 errors and green tests, and which errors were there before
it started.
\end{play}

\begin{play}{Review a project without changing anything}
\when You want to know where a project stands before you start on it.
\asks \ask{check the code standards, without changing anything}.
\happens \tool{sheriff_test} analyzes each component and gives the count per
file and per rule. Since you asked for nothing to change, the assistant stops
there.
\done You have the count and the code is as it was (\code{git status} clean).
\end{play}

\begin{play}{Fix what Sheriff fixes by itself}
\when A component with many formatting errors: imports, braces, names. It is
free and quick, so it is worth doing before anything else.
\asks \ask{fix what Sheriff can fix by itself}.
\happens \tool{sheriff_fix} applies every automatic repair Sheriff has, round
after round while the count drops, and reports the errors before and after and
which files it changed. It uses no model, so it \textbf{spends no tokens}. It
edits the files in place, with no commit.
\done The answer says \code{N error(s) before, M after}. Review the change with
\code{git diff}; to undo it, \code{git restore .}
\end{play}

\begin{play}{Read the rules before writing}
\when You are about to write something and want to know what Sheriff expects,
or an error makes no sense.
\asks \ask{what are the rules for javadoc?} or \ask{explain the SortedImport
rule}.
\happens \tool{sheriff_guidelines} searches the rule catalog. With a word, the
rules that match; with a rule's code, the whole rule and the code Sheriff
accepts for it. The first lookup on a machine takes a few seconds more: that is
when the catalog is extracted from the image.
\done You have the rule, its explanation and a correct example.
\end{play}

\begin{play}{One module of a project with several}
\when The repository has several modules (\code{api}, \code{web},
\code{core}\ldots) and you care about one.
\asks \ask{check the code standards of the api module}.
\happens The component is always the name of the module's folder, never a path
inside it; if the assistant passes a path, the tool answers with the right
component. The hooks only look at the components the turn changed.
\done As in play 1, for that module.
\end{play}

\begin{play}{Repair a whole component without steering each step}
\when A component with hundreds of errors, typically the first time an existing
project is analyzed, and you would rather it repaired itself while you do
something else.
\asks \ask{repair the whole component with sheriff\_autofix}.
\happens The assistant asks for your permission, because this tool
\textbf{commits and spends tokens}. It answers at once with a task id and works
in the background:
\begin{enumerate}
  \item It creates a branch of its own, \code{sheriff-agent/<date>}.
  \item It applies Sheriff's automatic repairs, in a commit of their own.
  \item It hands the model the files with errors in batches of five, one commit
        per pass, while the count keeps dropping.
  \item If the model touches a file outside its batch, it stops without
        committing and leaves that work in the tree for you to review.
  \item When it ends, however it ends, it runs the project's tests.
\end{enumerate}
To see how it is going, ask \ask{how is the task going?}: the assistant calls
\tool{sheriff_task}.
\done The task shows as finished. Review the branch (\code{git log},
\code{git diff main...}) and merge it if you like it.
\end{play}

\begin{warning}
Keep the session open while the repair works. In Codex, closing the session
cuts it short (section~\ref{sec:codex}).
\end{warning}

\begin{play}{Repair from the terminal, without an assistant}
\when You want to start the repair loop yourself, for example on a machine
with no assistant session open.
\asks From the module's folder:
\begin{command}[basicstyle=\ttfamily\footnotesize]{bash}
# The whole loop
java -jar ~/.local/share/sheriff-agent/sheriff-mcp.jar --agent
# Only analyze, or only the automatic repairs
java -jar ~/.local/share/sheriff-agent/sheriff-mcp.jar --agent --check-only
java -jar ~/.local/share/sheriff-agent/sheriff-mcp.jar --agent --sheriff-fix
# Search the rules
java -jar ~/.local/share/sheriff-agent/sheriff-mcp.jar --agent --rules javadoc
\end{command}
\happens The loop is the same as in the previous play: a branch of its own, a
commit per pass and the tests at the end. By default Claude Code answers, on
your subscription, with no API key. The \code{AI_BACKEND} variable picks
another: \code{codex_cli}, \code{antigravity_cli} (section~\ref{sec:antigravity}), \code{anthropic_api},
\code{openai_api} or \code{local} (a server compatible with OpenAI's API, such
as Ollama).
\done The final summary says how many errors are left and whether the tests
pass.
\end{play}

\begin{play}{Keep errors out in CI}
\when You want continuous integration to fail when the code does not comply.
\asks From the project's root:
\begin{command}{bash}
java -jar ~/.local/share/sheriff-agent/sheriff-mcp.jar --check
\end{command}
\happens It analyzes every component and ends with an exit code:

\begin{center}
\rowcolors{2}{pbcode}{white}
\begin{tabularx}{0.9\linewidth}{@{\hspace{6pt}}>{\ttfamily\bfseries}p{12mm}L}
\headrow \thead{Exit} & \thead{Meaning} \\
0 & Clean. \\
1 & There are errors; they are listed. \\
2 & It could not check: Docker is not available, or there is nothing of the
    profile's language to analyze. \\
\end{tabularx}
\end{center}

If Sheriff's image is not there, it pulls it before analyzing. In GitHub
Actions, for example, with no token and no account:

\begin{command}{GitHub Actions}
- uses: actions/setup-java@v5
  with: { distribution: temurin, java-version: '17' }
- run: |
    curl -fsSLO https://github.com/kaizten/sheriff-mcp-server/\
    releases/latest/download/sheriff-mcp.jar
- run: java -jar sheriff-mcp.jar --check
\end{command}
\done The CI step is green, or red with the list of errors.
\end{play}

\begin{play}{Keep errors out of the Maven build}
\when A Maven project where \code{mvn package} must fail on errors.
\asks Add the plugin to the \code{pom.xml}, with no configuration:
\begin{command}[escapeinside={(*}{*)}]{pom.xml}
<plugin>
  <groupId>com.kaizten</groupId>
  <artifactId>sheriff-maven-plugin</artifactId>
  <version>(*\pluginversion*)</version>
  <executions>
    <execution><goals><goal>check</goal></goals></execution>
  </executions>
</plugin>
\end{command}
\happens \code{sheriff:check} runs in the \code{package} phase, after the tests:
\code{mvn package}, \code{mvn verify} and \code{mvn install} fail on errors.
\code{mvn sheriff:fix} applies the automatic repairs and fails if anything is
left. Both leave Sheriff's JSON in \path{target/sheriff/}.
\done \code{mvn package} passes, or fails naming the errors.
\end{play}

\begin{note}
The Maven plugin is not in the release: it is installed into \path{~/.m2} by
building the tools' repository (\code{mvn install} in a clone). For CI,
\code{--check} needs nothing but the jar.
\end{note}

% ---------------------------------------------------------------------------
\section{The hooks: what you see when they stop you}
\label{sec:hooks}

The assistant uses the MCP tools when it sees fit. The hooks are different:
Claude Code and Codex run them at set moments, and they are what makes sure the
work is not called done with errors left. The installer turns them on in both,
for all your projects; in Codex, once you approve them (section~\ref{sec:codex}).
Antigravity has no Sheriff hooks.

\begin{center}
\rowcolors{2}{pbcode}{white}
\begin{tabularx}{\linewidth}{@{\hspace{6pt}}K{24mm}>{\raggedright\arraybackslash}p{35mm}L}
\headrow \thead{Hook} & \thead{When it acts} & \thead{What it does} \\
Gate & Before the session's first edit to a component &
  If the component already has errors, it blocks that edit once, with the
  list, so that they are cleared first. The edits after it go through, and an
  edit to a file that has errors is never blocked: it is the fix. Claude
  Code only: Codex edits with patches, and the Stop hook catches the same
  after. \\
Stop & When the assistant is about to finish &
  If what the turn changed has errors, or new methods with no test (in Java),
  it sends the turn back with the list. \\
Turn start & When your message arrives &
  Records how the components looked, so that the Stop hook looks only at what
  changes in this turn, not at earlier uncommitted work. \\
\end{tabularx}
\end{center}

When a hook stops the assistant, you see a message like these in the
conversation, followed by the same kind of \code{Next step:} the tools give:

\begin{answer}[Gate]
Blocked: basket already has 4 Sheriff error(s), from before this change, and basket/src/main/java/demo/Discount.java
has none of them. They are fixed first, in every file of the component, before the
new code. Edits to the files that have them go through; make this one again once
they are fixed, or now if it is part of fixing them.
\end{answer}

\begin{answer}[Stop]
Not done: what this turn changed has Sheriff errors, those that were already there
included. All of them are fixed before finishing; only a false positive of the checker
may stay, named, with the reason.
\end{answer}

\subsection{How they behave}

\begin{itemize}
  \item \textbf{They never block you over a failure of their own.} Without Java
        or Docker, the work goes on and one warning line is left.
  \item \textbf{Where there is nothing Sheriff analyzes}, they let everything
        through in a tenth of a second.
  \item \textbf{They do not loop.} The Stop hook sends the turn back at most
        five times in a row, and if the assistant cannot change anything, it
        lets the turn end.
  \item \textbf{They stand aside for the repair agent}, which runs its own
        checks.
\end{itemize}

\subsection{Require the tests at the end too}

By default the Stop hook requires 0 errors but does not run the tests, which can
take minutes. To require them too, export this variable before opening Claude
Code. With Codex it cannot be done: Codex does not pass your terminal's
variables to its hooks.

\begin{command}{bash}
export SHERIFF_STOP_RUNS_TESTS=1
claude
\end{command}

\subsection{See them, remove them, or keep them to one project}

\begin{itemize}
  \item \textbf{See the active hooks:} \code{/hooks}, inside Claude Code or
        Codex.
  \item \textbf{Remove them:}
\begin{command}{bash}
java -jar ~/.local/share/sheriff-agent/sheriff-mcp.jar --uninstall-hooks --user
\end{command}
        In Codex, \code{--uninstall-codex} removes them together with the server
        and the order of work (section~\ref{sec:uninstall}).
  \item \textbf{Have them in one project only}, in Claude Code: install with \code{--no-hooks}
        and, from that project's root, run:
\begin{command}{bash}
java -jar ~/.local/share/sheriff-agent/sheriff-mcp.jar --install-hooks
\end{command}
        They are written into its \path{.claude/settings.json}. That file names
        the jar's path: commit it only if the whole team installs to the same
        path.
\end{itemize}

% ---------------------------------------------------------------------------
\section{Configuring your project}
\label{sec:configure}

Nothing needs configuring. This section is for when you want something other
than what the tools work out by themselves.

\subsection{What is analyzed}

\begin{center}
\rowcolors{2}{pbcode}{white}
\begin{tabularx}{\linewidth}{@{\hspace{6pt}}>{\bfseries}p{24mm}>{\ttfamily}p{24mm}>{\ttfamily}p{30mm}L}
\headrow \thead{Language} & \thead{Files} & \thead{Base profile} & \thead{Architectures} \\
Java & .java & JAVA & \code{JAVA_HEXAGONAL}, \code{JAVA_DDD} \\
TypeScript & .ts, .tsx & TYPESCRIPT & \code{TYPESCRIPT_HEXAGONAL} \\
Vue & .vue & VUEJS & \\
Perl & .pl, .pm & PERL\_FORMAT & \\
Python & .py & PYTHON & \\
\end{tabularx}
\end{center}

A component in any other language is not passed as good: the tools say there
was nothing to analyze. If the project builds with Maven, Gradle or an npm
\code{test} script, its tests are run; if not, it is analyzed all the same and
the tools say that no tests were run.

\subsection{Declare the architecture}

With nothing declared, a Java project is checked under \code{JAVA}: style and
documentation. If it follows a hexagonal or DDD architecture, declare it once
and every tool (server, hooks, \code{--check} and Maven plugin) takes it into
account. In the \code{pom.xml}, which the modules inherit:

\begin{command}{pom.xml}
<properties>
  <sheriff.profile>JAVA_HEXAGONAL</sheriff.profile>
</properties>
\end{command}

Or, in a project without Maven, in a \code{.sheriff.properties} file beside its
sources:

\begin{command}{.sheriff.properties}
profile=TYPESCRIPT_HEXAGONAL
\end{command}

The architecture's rules \textbf{are added} to the language's, never put in
their place. Several can be listed, separated by commas
(\code{JAVA_HEXAGONAL,JAVA_DDD}); each one adds one more analysis.

\subsection{Every profile}

Sheriff has more profiles than the ones above. This is the full list of its
image of 8 October 2026, with the base profile in bold. Any of them is declared
like an architecture, and the tools always add its language's base profile:
declaring \code{JAVA_HEXAGONAL_REST} runs \code{JAVA} and
\code{JAVA_HEXAGONAL_REST}.

\begin{center}
\rowcolors{2}{pbcode}{white}
\begin{tabularx}{\linewidth}{@{\hspace{6pt}}>{\bfseries}p{20mm}L}
\headrow \thead{Language} & \thead{Profiles} \\
Java & \textbf{\code{JAVA}}, \code{JAVA_HEXAGONAL}, \code{JAVA_HEXAGONAL_DOMAIN}, \code{JAVA_HEXAGONAL_APPLICATION}, \code{JAVA_HEXAGONAL_REST}, \code{JAVA_HEXAGONAL_MONGODB}, \code{JAVA_HEXAGONAL_TIMESCALE}, \code{JAVA_DDD}, \code{JAVA_DDD_ENTITY}, \code{JAVA_DDD_VALUE_OBJECT}, \code{JAVA_DDD_ENUMERATE}, \code{JAVA_ENTITY}, \code{JAVA_VALUE_OBJECT}, \code{JAVA_ENUMERATE}, \code{JAVA_USE_CASE}, \code{JAVA_POM}, \code{JAVA_COMPILATION} \\
TypeScript & \textbf{\code{TYPESCRIPT}}, \code{TYPESCRIPT_HEXAGONAL}, \code{TYPESCRIPT_HEXAGONAL_DOMAIN}, \code{TYPESCRIPT_HEXAGONAL_APPLICATION}, \code{TYPESCRIPT_HEXAGONAL_HTTP}, \code{TYPESCRIPT_HEXAGONAL_VUEJS}, \code{TYPESCRIPT_ENTITY}, \code{TYPESCRIPT_VALUE_OBJECT}, \code{TYPESCRIPT_ENUMERATE}, \code{TYPESCRIPT_FILENAME}, \code{TYPESCRIPT_LOCALE}, \code{TYPESCRIPT_PACKAGE}, \code{TYPESCRIPT_COMPILATION} \\
Vue & \textbf{\code{VUEJS}}, \code{VUEJS_COMPONENT}, \code{VUEJS_VIEW}, \code{VUEJS_FILENAME} \\
Perl & \textbf{\code{PERL_FORMAT}} \\
Python & \textbf{\code{PYTHON}} \\
\end{tabularx}
\end{center}

To see the list of the Sheriff version you have, ask it for a profile that does
not exist. After an error line, the longest \code{Possible Values} line is the
profiles:

\begin{command}{bash}
docker run --rm kaizten/sheriff:latest test -t LIST
\end{command}

What each profile checks is Kaizten's to define. The tools know the exact rules
of 12 of them, the ones Sheriff's own export covers.

To try a profile for one request, without declaring it, name it:
\ask{check the api module with JAVA\_DDD}. The tools use it for that request,
and the hooks keep to the profile the project declares.

\subsection{No rule can be switched off}

What the profile holds is what is checked. The one exception is a known false
positive of Sheriff's, \code{FilenamePascalCase} on \code{package-info.java} and
\code{module-info.java}, names Java itself requires, which the tools do not
show. If the assistant finds another false positive, it may only leave it by
saying which one it is and why.

\subsection{Share the setup with the team}

Each machine has the server registered by the installer. In Claude Code, if
you would rather the project declared it, commit a \code{.mcp.json} at its root
with the jar's path. Use one or the other: with both, Claude Code warns of two \code{sheriff}
servers in different scopes.

\begin{command}[basicstyle=\ttfamily\footnotesize]{.mcp.json}
{
  "mcpServers": {
    "sheriff": {
      "command": "java",
      "args": ["-jar", "/home/<your-user>/.local/share/sheriff-agent/sheriff-mcp.jar"]
    }
  }
}
\end{command}

\subsection{Useful variables}

With Claude Code, exporting them before you open it is enough. With Codex, they
do not arrive from your terminal (section~\ref{sec:codex}).

\begin{center}
\rowcolors{2}{pbcode}{white}
\begin{tabularx}{\linewidth}{@{\hspace{6pt}}>{\ttfamily\small}p{50mm}L}
\headrow \thead{Variable} & \thead{What for} \\
SHERIFF\_PULL=always & Pull a newer Sheriff image at startup (with
  \code{never}, Docker Hub is not even asked). \\
SHERIFF\_IMAGE & Pin the image by its digest
  (\code{kaizten/sheriff@sha256:...}), so that the rules do not change until
  you decide. \\
SHERIFF\_TIMEOUT & Seconds per analysis; 300 by default. \\
SHERIFF\_STOP\_RUNS\_TESTS=1 & The Stop hook requires green tests too. \\
SHERIFF\_STOP\_MAX\_BLOCKS & How many times in a row it sends the turn back; 5
  by default. \\
VERIFICATION\_TEST\_CMD & The test command, when it is not the one the build
  tool gives. \\
AI\_BACKEND & Who answers in the terminal's repair loop. \\
\end{tabularx}
\end{center}

% ---------------------------------------------------------------------------
\section{Codex}
\label{sec:codex}

When \code{codex} is installed, the installer sets it up as it does Claude
Code, for every project. It writes three things into \path{~/.codex/}, leaving
the rest of those files as they were:

\begin{itemize}
  \item in \code{config.toml}, the server, with the time limit it needs (Codex
        cuts a call at 60 seconds) and \tool{sheriff_fix} approved;
  \item in \code{AGENTS.md}, the order of work, which Codex reads in every
        project;
  \item in \code{hooks.json}, the Stop and turn-start hooks. Not the gate,
        because Codex edits with patches; the Stop hook catches the same errors
        after the edit.
\end{itemize}

\subsection{What is left to you}

\textbf{Approve the hooks once.} Codex asks before it runs a new hook: you see
\code{Hooks need review} in the first session, or you can do it with
\code{/hooks}. Until you approve them, the Stop hook does not run.

\begin{warning}
Without the Stop hook, Codex can call the work done with errors left. On
PetClinic, asked for a single method, it stopped halfway twice, with 133 and
120 errors to go; the Stop hook is what kept it going down to 0.
\end{warning}

\subsection{Four things worth knowing}

\begin{itemize}
  \item \textbf{Open Codex in the project's folder.} Codex does not say where the
        project is, so it is worked out from the folder it runs in.
  \item \textbf{\code{codex exec} does not write unless told to.} To let it
        edit: \code{codex exec -s workspace-write "Add a method size() to Basket"}.
  \item \textbf{\tool{sheriff_autofix} ends with the session.} Codex stops its
        servers when a session ends, and the repair loop with them, which can
        leave the repository on the loop's branch with a pass uncommitted. For
        long repairs, use the terminal:
\begin{command}[basicstyle=\ttfamily\footnotesize]{bash}
AI_BACKEND=codex_cli java -jar ~/.local/share/sheriff-agent/sheriff-mcp.jar --agent
\end{command}
  \item \textbf{Your terminal's variables do not arrive} at the server or the
        hooks. Use the installer's options (\code{--pull-always},
        \code{--fail-fast}) or write them into the server's entry in
        \code{config.toml}: \texttt{env = \{ SHERIFF\_TIMEOUT = "600" \}}.
\end{itemize}

% ---------------------------------------------------------------------------
\section{Antigravity}
\label{sec:antigravity}

Google's Antigravity can work with Sheriff in two ways: as the assistant, with
the MCP server, as Claude Code and Codex do, and as one of the repair loop's
models. Both go through its command line, \code{agy}.

\subsection{As the assistant}

When \code{agy} is installed, the installer sets it up, for every project. It
writes two things, leaving the rest of those files as they were:

\begin{itemize}
  \item in \path{~/.gemini/config/mcp_config.json}, the server, the same entry
        \code{agy mcp add} would register;
  \item in \path{~/.gemini/antigravity-cli/settings.json}, permission for the
        four tools that neither commit nor spend tokens to run without asking,
        as in Claude Code: \code{mcp(sheriff/sheriff_test)} and the other
        three. \tool{sheriff_autofix} always asks.
\end{itemize}

If you install Antigravity later, run the install line again, or:

\begin{command}[basicstyle=\ttfamily\footnotesize]{bash}
java -jar ~/.local/share/sheriff-agent/sheriff-mcp.jar --install-antigravity
\end{command}

Antigravity reads the server's order of work, so it follows the same cycle as
Claude Code (section~\ref{sec:first}). What it does not have is the hooks: nothing stops an Antigravity turn from
ending with errors left; it depends on the model following the steps.

\begin{note}
Tried on 7 October on Linux, with \code{agy} 1.3.1 and the \code{basket}
project of section~\ref{sec:first}. Asked to \ask{fix this project's code
standards errors}, it called \tool{sheriff_test} and \tool{sheriff_fix}, fixed
by hand what was left, checked again, ran \code{mvn test} and ended with 0
errors.
\end{note}

\subsection{In the repair loop}

With \code{agy} logged in, from the module's folder:

\begin{command}[basicstyle=\ttfamily\footnotesize]{bash}
AI_BACKEND=antigravity_cli java -jar ~/.local/share/sheriff-agent/sheriff-mcp.jar --agent
\end{command}

The loop is the same as with Claude: a branch of its own, Sheriff's automatic
repairs first, one commit per pass and the project's tests at the end.
\needspace{6\baselineskip}
To pick the model (\code{agy models} lists them):

\begin{command}{bash}
export ANTIGRAVITY_COMMAND="agy --mode accept-edits --model claude-sonnet-5-5-high"
\end{command}

In the loop, Antigravity is confined:

\begin{itemize}
  \item \textbf{It edits only inside the component.} Antigravity itself refuses
        anything outside it.
  \item \textbf{It can run two commands only}: the project's tests and
        \code{git mv}. If it tries another, or a URL, Antigravity refuses it and
        the pass ends; the loop records it and goes on.
  \item \textbf{Your Antigravity settings are left alone.} Each pass works with
        a temporary copy of them, deleted when it ends.
  \item \textbf{Each pass has 600 seconds.} The
        \code{SHERIFF_AGENT_ANTIGRAVITY_TIMEOUT} variable changes that limit.
  \item \textbf{It has not been tried on Windows}, so there it runs no command,
        not even the tests.
\end{itemize}

\begin{note}
Tried on 7 October on Spring PetClinic: from 221 errors to 0 in five passes,
with the 81 tests and the project's formatter green.
\end{note}

% ---------------------------------------------------------------------------
\section{Troubleshooting}
\label{sec:troubleshooting}

\newcommand\problem[2]{\par\addvspace{0.5\baselineskip}\noindent\textbf{\color{pbbrand}#1}\par\nopagebreak\vspace{-0.45\baselineskip}\noindent\hangindent=1.2em\hangafter=0 #2\par}

\subsection{When installing}

These are the installer's own messages. On the first four it stops; on the rest
it goes on and warns you.

\problem{\code{No Java 17 or newer was found}}{There is no Java 17 or newer in
\code{JAVA_HOME}, on the \code{PATH} or in the folders JDKs are installed in.
Install one (for example, Temurin from \path{https://adoptium.net}; on Windows,
\code{winget install EclipseAdoptium.Temurin.21.JDK}) and run the line again.
It need not be on the \code{PATH}.}
\problem{\code{... could not be downloaded}}{There is no network, or there is
a proxy \code{curl} or PowerShell does not know about. With
\code{--release=TAG} (\code{-Tag}), there is no release by that name. Run the
line again once there is network.}
\problem{\code{the jar downloaded does not match its checksum}}{The download
was cut short. Run the line again.}
\problem{\code{The MCP server does not start: FAIL}}{The installer started the
server with the Java it found, as your assistant will, and it did not answer.
What the server said follows; with that fixed, run the line again.}
\problem{\code{Docker is not installed. The install goes on}\ldots}{The tools
are installed, but they analyze nothing until there is Docker. Install it, or
run the line again with \code{--install-docker}.}
\problem{\code{Docker is not running. The install goes on}\ldots}{Start
Docker; the server pulls the image when it starts. On Linux, if Docker is in
fact running, most likely your user cannot use it without \code{sudo}: add
yourself to the group with \code{sudo usermod -aG docker $USER}, log out of the
system and back in.}
\problem{\code{Sheriff's image could not be pulled now}}{There is no network
or, on Windows, Docker Desktop is set to Windows containers
(\code{no matching manifest for windows}): switch it to Linux containers, its
default. The server pulls the image when it starts.}
\problem{\code{claude}, \code{codex} or \code{agy is not on the PATH}}{That
assistant was not installed, so the installer did not set it up. Once you
install it, run the line again, or the command it printed.}
\problem{\code{No assistant was found}}{None of Claude Code, Codex or
Antigravity is installed, so the installer wrote nothing into any assistant's
configuration: Sheriff is there only as the CI check (\code{--check}) and the
Maven plugin. Install one and run the line again.}

\subsection{When a session opens}

\problem{\code{Failed to connect} in \code{claude mcp list}, or \code{connection closed}}{The
jar is not where the registration says, or the Java registered is gone: when
Java is updated on Windows, the JDK's folder changes its name. Run the install
line again, which looks for Java again and checks that the server starts.}
\problem{On Windows, the hooks fail with a \code{non-blocking} error and
\code{java: command not found}}{An earlier installer wrote them, with
\code{java} and no path, and Claude Code runs the hooks in Git Bash, where Java
may not be on the \code{PATH}. Run the install line again: it now writes the
full path of \code{java.exe}.}
\problem{Claude Code asks for permission every time it uses \tool{sheriff_fix}}{The
permissions were not written: install again. \tool{sheriff_autofix} always
asks, on purpose.}
\problem{The hooks do nothing}{The session has not reloaded its settings.
Restart it and check them with \code{/hooks}.}

\subsection{When analyzing}

\problem{\code{Sheriff's image ... is being downloaded}}{It is the first use on
this machine: wait a few minutes.}
\problem{\code{The analyzer could not run}}{Docker has stopped. Start it and ask
for the check again.}
\problem{\code{Nothing was analyzed}}{The component holds no sources of the
profile's language. Name the component (\ask{check the api module}) or check
the profile the project declares.}
\problem{\code{No rule catalog is available}}{There was no Docker when the
server started. Start Docker: the next rule lookup tries again, at most a
minute later.}
\problem{\code{a newer kaizten/sheriff:latest has been published}}{There are new
rules. Run the install line again to pull them, or install with
\code{--pull-always}.}

\subsection{With Codex}

\problem{\code{timed out awaiting tools/call}}{The server's entry has no time
limit. Install again, or run the jar with \code{--install-codex}.}
\problem{\code{writing is blocked by read-only sandbox}}{\code{codex exec} is
read-only by default: add \code{-s workspace-write}.}
\problem{The Stop hook does not act}{The hook has not been approved. Open an
interactive session and approve it with \code{/hooks}.}

\subsection{How to report a problem}

Include three things: the version (the jar with \code{--version}), the task id
the answer shows (\code{Task 4047f176.}) and, if you can, the file
\path{~/.local/state/sheriff-mcp/tasks/<id>/task.json}, which records the jar's
version and the exact Sheriff image the analysis ran with. The newest 50 tasks
are kept.

% ---------------------------------------------------------------------------
\section{Frequently asked questions}
\label{sec:faq}

\newcommand\faq[2]{\par\addvspace{0.6\baselineskip}\noindent{\sffamily\bfseries\color{pbbrand}#1}\par\nopagebreak\vspace{-0.45\baselineskip}#2\par}

\faq{Does this spend tokens?}{Only \tool{sheriff_autofix} and the terminal's
loop, which call a model. Analyzing, reading the rules and the automatic
repairs of \tool{sheriff_fix} use no model. The work your assistant does by
hand on the errors left is part of your normal session.}

\faq{Is anything copied into my project?}{No. Sheriff leaves a few temporary
files while it analyzes, and the tools delete them after. The only things that
change are the files \tool{sheriff_fix} or the assistant repair, and you see
those with \code{git diff}.}

\faq{Why does it fix errors that were not mine?}{Because the component is
delivered clean as a whole, and leaving the \emph{old} errors is the most
common way for them never to be fixed. When you do not want it, say so in the
request.}

\faq{Why is the first time so slow?}{Sheriff's image takes about 4\,GB and the
rule catalog is extracted from it (about 6 seconds). After that, an analysis
takes a few seconds.}

\faq{Do new rules arrive by themselves?}{No, on purpose: rules that change in
the middle of a piece of work confuse. When there is a newer image, it is
reported once per session. To pull it, install again; to always pull it,
install with \code{--pull-always}.}

\faq{What if Docker is not running?}{There is no analysis. The tools say so,
the assistant goes on with the work and runs the tests, and the hooks let you
through.}

\faq{Does it work with other assistants?}{Any MCP client can use the server.
For one that does not read MCP instructions, the jar with
\code{--instructions} prints the order of work, and \code{--call sheriff_test}
runs a tool from the terminal with the same answer the server gives.}

% ---------------------------------------------------------------------------
\section{Update and uninstall}
\label{sec:uninstall}

\textbf{To update}, run the install line again: it downloads the latest release
and writes the configuration again. Then restart the open sessions. Run it
again too when you update or change Java: the assistants call it by its full
path, and the old Java's may disappear.

In the uninstall commands, if \code{java} is not on the \code{PATH}, put in its
place the path the installer showed on its first line
(\code{Using Java ... at ...}).

\textbf{To uninstall}, in this order, as the first steps use the jar. On
Linux and macOS:

\begin{command}{bash}
java -jar ~/.local/share/sheriff-agent/sheriff-mcp.jar --uninstall-hooks --user
# If you use Codex
java -jar ~/.local/share/sheriff-agent/sheriff-mcp.jar --uninstall-codex
# If you use Antigravity
java -jar ~/.local/share/sheriff-agent/sheriff-mcp.jar --uninstall-antigravity
claude mcp remove --scope user sheriff
rm -rf ~/.local/share/sheriff-agent ~/.cache/sheriff-mcp ~/.local/state/sheriff-mcp
# Optional: frees the image's 4 GB
docker rmi kaizten/sheriff:latest
\end{command}

On Windows, in PowerShell:

\begin{command}[basicstyle=\ttfamily\footnotesize]{PowerShell}
$jar = "$HOME\.local\share\sheriff-agent\sheriff-mcp.jar"
java -jar $jar --uninstall-hooks --user
# If you use Codex
java -jar $jar --uninstall-codex
# If you use Antigravity
java -jar $jar --uninstall-antigravity
claude mcp remove --scope user sheriff
Remove-Item -Recurse -Force -ErrorAction SilentlyContinue "$HOME\.local\share\sheriff-agent"
Remove-Item -Recurse -Force -ErrorAction SilentlyContinue "$HOME\.cache\sheriff-mcp"
Remove-Item -Recurse -Force -ErrorAction SilentlyContinue "$HOME\.local\state\sheriff-mcp"
# Optional: frees the image's 4 GB
docker rmi kaizten/sheriff:latest
\end{command}

% ---------------------------------------------------------------------------
\clearpage
\section{Quick reference card}
\label{sec:card}

\begin{minipage}[t]{0.48\linewidth}\raggedright
\cardhead{What to ask}
\begin{itemize}[leftmargin=1.1em]
  \item \ask{check the code standards}
  \item \ask{fix what Sheriff can fix by itself}
  \item \ask{what are the rules for javadoc?}
  \item \ask{check the code standards of the api module}
  \item \ask{repair the whole component with sheriff\_autofix}
  \item \ask{how is the task going?}
  \item \ask{\ldots{} leave the errors that were already there}
\end{itemize}

\cardhead{The five tools}
\rowcolors{2}{pbcode}{white}
\begin{tabularx}{\linewidth}{@{\hspace{4pt}}>{\ttfamily\small}p{33mm}>{\small\raggedright\arraybackslash}X}
\headrow \thead{Tool} & \thead{What it does} \\
sheriff\_test & Analyzes: count and next step. \\
sheriff\_fix & Automatic repairs; lists what is left. \\
sheriff\_guidelines & The rules and the code accepted. \\
sheriff\_autofix & Repairs it all on a branch. \textbf{Spends tokens.} \\
sheriff\_task & A task's state. \\
\end{tabularx}

\cardhead{Exit codes of \code{--check}}
\rowcolors{2}{pbcode}{white}
\begin{tabularx}{\linewidth}{@{\hspace{4pt}}>{\ttfamily\bfseries\small}p{8mm}>{\small}X}
\headrow \thead{} & \thead{Meaning} \\
0 & Clean. \\
1 & There are errors. \\
2 & Could not check. \\
\end{tabularx}
\end{minipage}%
\hfill
\begin{minipage}[t]{0.48\linewidth}\raggedright
\cardhead{Install or update}
{\ttfamily\footnotesize curl -fsSL https://github.com/\textbackslash\\
kaizten/sheriff-mcp-server/\textbackslash\\
releases/latest/download/\textbackslash\\
install.sh | bash\par}
{\footnotesize On Windows, the line in section~\ref{sec:install}.\par}

\cardhead{Commands}
{\ttfamily\footnotesize
\begin{tabular}{@{}l@{}}
\textcolor{pbmuted}{J=\textasciitilde/.local/share/sheriff-agent/sheriff-mcp.jar}\\[3pt]
java -jar \$J -{}-version\\
java -jar \$J -{}-check\\
java -jar \$J -{}-call sheriff\_test\\
java -jar \$J -{}-agent\\
AI\_BACKEND=antigravity\_cli java -jar \$J -{}-agent\\
java -jar \$J -{}-install-hooks [-{}-user]\\
java -jar \$J -{}-uninstall-hooks [-{}-user]\\
java -jar \$J -{}-install-codex\\
java -jar \$J -{}-install-antigravity\\
\end{tabular}\par}

\cardhead{Where things are}
\rowcolors{2}{pbcode}{white}
\begin{tabularx}{\linewidth}{@{\hspace{4pt}}>{\ttfamily\scriptsize}p{43mm}>{\small\raggedright\arraybackslash}X}
\headrow \thead{Path} & \thead{What} \\
\textasciitilde/.local/share/sheriff-agent & The jar. \\
\textasciitilde/.cache/sheriff-mcp & The rule catalog. \\
\textasciitilde/.local/state/sheriff-mcp & The tasks. \\
\textasciitilde/.claude/settings.json & Hooks and permissions. \\
\textasciitilde/.codex/ & Codex's configuration. \\
\textasciitilde/.gemini/config/ & The server in Antigravity. \\
\end{tabularx}

\cardhead{Done means}
\begin{itemize}[leftmargin=1.1em]
  \item \textbf{0 errors} in the component,
  \item \textbf{green tests},
  \item and a test for every new method.
\end{itemize}
\end{minipage}

\end{document}
